\chapter*{Conclusioni} %Se si cambia il Titolo cambiare anche la riga successiva così che appaia corretto nell'indice
\addcontentsline{toc}{chapter}{Conclusioni} %Per far apparire Conclusioni nell'indice (Il nome deve rispecchiare quello del chapter)

Nessuna delle protezioni di sicurezza attive di default sul binario (Codice~\ref{lst:checksec_output}, Tabella~\ref{tab:ambiente_riferimento}) impedisce di causare leak di informazioni e di eseguire scritture in memoria. Delle due protezioni pensate specificamente per questa vulnerabilità (Capitolo~\ref{chapter:mitigations}), quella a runtime (Sezione~\ref{sec:fortify_source}) blocca la scrittura arbitraria ma lascia comunque ottenibile il leak di informazioni, anche se con un payload più lungo e meno mirato; solo quella sul codice sorgente (Sezione~\ref{sec:prevention}), abbinata a \texttt{-Werror}, impedisce che il pattern vulnerabile venga anche solo compilato. Ne emerge una gerarchia netta: le protezioni generiche sono inefficaci contro questa vulnerabilità, quelle dedicate a runtime ne riducono il danno senza eliminarla e solo l'intervento sul codice sorgente la previene alla radice.

Questo risultato non è solo argomentato ma dimostrato: la Sezione~\ref{sec:fortify_verify} riprende lo stesso exploit del Capitolo~\ref{chapter:exploit} contro lo stesso programma ricompilato con la protezione attiva. È questa continuità a rendere il risultato più solido di una trattazione puramente teorica.

Il lavoro resta circoscritto all'ambiente di riferimento (Tabella~\ref{tab:ambiente_riferimento}): un'estensione naturale sarebbe verificare se lo stesso meccanismo vale su un'altra architettura. Un esempio è AArch64, la cui calling convention riserva otto registri general-purpose agli argomenti anziché sei, con una register save area equivalente definita esplicitamente per le funzioni variadiche~\cite{aapcs64}: il numero di argomenti fantasma prima di esaurire i registri diventerebbe sette anziché cinque, ma il meccanismo della Sottosezione~\ref{subsec:funvar_cc} resterebbe identico.
